\bookchapter{Capítulo 50: Kotlin em ritmo acelerado — a primeira linguagem do desenvolvimento Android}

\chapterintro{Neste capítulo vamos estudar a linguagem de verdade. Sem medo — o Kotlin é famoso por ser “amigável para iniciantes” e você não precisa decorar a gramática: dominando 20\% da sintaxe, você já entende e escreve 80\% do código do dia a dia. O resto você consulta enquanto desenvolve o projeto. Lembre do nosso princípio de estudo: entender > decorar, saber usar > ser perfeito.}

\booksection{Variáveis e tipos: dando nome aos dados}

Um programa é “dados + lógica”. Variável é dar um nome aos dados para guardá-los. Em Kotlin há duas formas de declarar:

\begin{lstlisting}[language=Kotlin]
val name = João        // val: variável somente leitura (o padrão recomendado)
var age = 18             // var: variável mutável (pode mudar depois)

age = 19                 // <i class="fas fa-check-circle mr-1"></i> var pode mudar
// name = "Maria"         // <i class="fas fa-times-circle mr-1"></i> val não pode mudar: erro de compilação

val score: Int = 100     // também é possível declarar o tipo explicitamente
val pi = 3.14            // se omitir, é inferido automaticamente
\end{lstlisting}

Os tipos mais usados são poucos; memorize com uma tabela:

\begin{center}
\begin{tabularx}{\linewidth}{|X|X|X|}
\hline
\textbf{Tipo} & \textbf{O que guarda} & \textbf{Exemplo} \tabularnewline
\hline
Int & Números inteiros & 42 \tabularnewline
Double & Números decimais & 3.14 \tabularnewline
String & Texto & "Olá" \tabularnewline
Boolean & Verdadeiro/falso & true / false \tabularnewline
List< T> & Um conjunto de dados & listOf(1, 2, 3) \tabularnewline
\hline
\end{tabularx}
\end{center}

\begin{center}\small\color{gray}Tipos básicos mais usados do Kotlin\end{center}

\begin{tcbinfo}{Dica: o princípio do val primeiro}
A regra de ferro da comunidade Kotlin: por padrão use val e só use var quando realmente for necessário. Se usar var onde daria val, o IDE avisa com texto amarelo. Esse hábito elimina uma grande família de bugs de “variável alterada sem querer” — quanto menos o código muda, mais fácil raciocinar sobre ele.

\end{tcbinfo}

\booksection{Funções: empacotando a lógica em “botões”}

Uma função é um trecho de lógica com nome, como um botão do controle remoto — aperta (chama), e uma sequência de ações acontece. Em Kotlin, declaramos com fun:

\begin{lstlisting}[language=Kotlin]
fun greet(name: String): String {
    return "Olá, $name!"     // o símbolo $ insere a variável na string
}

println(greet(João))        // imprime: Olá, João!

// funções de uma linha podem ser escritas como expressões
fun add(a: Int, b: Int) = a + b

println(add(3, 5))            // imprime: 8
\end{lstlisting}

Alguns pontos-chave: os parâmetros ficam entre parênteses (nome primeiro, tipo depois, separados por dois-pontos); o tipo de retorno vem depois dos parênteses; e o template de string \$nomeDaVariavel é a conveniência marca registrada do Kotlin — sem precisar concatenar “string + variável” como no Java.

\booksection{Condicionais e laços: ensinando o programa a “pensar”}

A inteligência de um programa vem de duas coisas: saber decidir (if/when) e saber repetir (for/while).

\begin{lstlisting}[language=Kotlin]
// Decisão: if / else
val age = 17
if (age >= 18) {
    println("adulto")
} else {
    println("menor de idade")
}

// Multiramo: when (cem vezes mais prático que switch)
val score = 85
val level = when {
    score >= 90 -> "excelente"
    score >= 60 -> "na média"
    else -> "Continue assim!"
}

// Repetição: um loop for
val fruits = listOf("maçã", "banana", "laranja")
for (fruit in fruits) {
    println("Eu gosto de $fruit")
}
\end{lstlisting}

O when é a forma mais usada de múltiplas escolhas em Kotlin; no desenvolvimento Android você vai usá-lo todos os dias (por exemplo, para mostrar telas diferentes conforme o estado).

\booksection{Segurança contra nulos: o “amuleto de proteção” do Kotlin}

O bug mais famoso da era Java chama-se NullPointerException (exceção de ponteiro nulo) — o programa acessa algo “que não existe” e trava na hora. Nas estatísticas, ele lidera constantemente o ranking de travamentos do Android. O Kotlin resolveu isso no nível da linguagem: por padrão, valores nulos não são permitidos.

\begin{lstlisting}[language=Kotlin]
var name: String = "João"   // tipo comum: nunca pode ser null
// name = null                // <i class="fas fa-times-circle mr-1"></i> erro de compilação — te barra na hora

var nickname: String? = null // adicione um ponto de interrogação: null permitido
nickname = "Zé"

// Com tipos anuláveis, é preciso tratar com segurança:
println(nickname?.length)    // chamada segura: retorna null em vez de travar
val len = nickname?.length ?: 0   // operador Elvis: valor padrão quando null
println(len)
\end{lstlisting}

\begin{tcbwarn}{Atenção: não confunda ? e ?}
String? significa “pode ser nulo”; a ?: b significa “se a for nulo, use b” (o símbolo lembra o cabelo do Elvis, por isso se chama operador Elvis). No começo, quando o compilador reclamar de “erro relacionado a null”, quase certamente foi você quem esqueceu de tratar um tipo anulável — é justamente a linguagem te protegendo.

\end{tcbwarn}

\booksection{Classes e objetos: o molde que fabrica “peças”}

A classe é o “molde” do objeto, e o objeto é o “produto pronto” da classe. Tudo no Android — telas, botões, dados — é um objeto. O Kotlin tem uma joia chamada data class (classe de dados), feita sob medida para “guardar dados”: uma linha vale dezenas em Java:

\begin{lstlisting}[language=Kotlin]
// Classe comum: um molde
class Dog(val name: String) {
    fun bark() = println("$name: Au au!")
}

val dog = Dog("Rex")      // cria um objeto
dog.bark()                 // Rex: Au au!

// data class: uma caixa para dados (com equals/copy automáticos)
data class User(val id: Int, val name: String, val email: String)

val me = User(1, João, "zhangsan@example.com")
val twin = me.copy(name = "Joãozinho")   // copia e muda um campo
\end{lstlisting}

Quando você for fazer o seu aplicativo de controle de gastos, cada lançamento será uma data class: data class Record(val amount: Double, val category: String, val date: String). Viu? Você já consegue escrever a estrutura de código de um projeto real.

\booksection{Lambdas e funções de alta ordem: o “açúcar sintático” do Kotlin}

No código Android você verá todos os dias funções “sem nome” como \{ ... \}; elas se chamam lambdas (funções anônimas). O onClick = \{ count++ \} de um Button é um lambda — empacota “o que fazer ao clicar” num pacote só e entrega ao botão.

\begin{lstlisting}[language=Kotlin]
// Jeito comum: definir uma função com nome e passar
fun double(x: Int) = x * 2
val result = listOf(1, 2, 3).map(::double)

// Com lambda: definida na hora, leve e direta
val result2 = listOf(1, 2, 3).map { it * 2 }   // it = elemento atual

// Combinação filtro + transformação (comum demais em projetos reais)
val bills = listOf(12.5, 88.0, 5.0, 200.0)
val bigBills = bills.filter { it > 50 }.sorted()
// [88.0, 200.0] — "contas acima de 50, ordenadas"
\end{lstlisting}

Guarde alguns “combos de coleção” muito usados: filter (filtrar), map (transformar), sorted (ordenar) e sum (somar). Nos próximos tempos, 90\% do processamento de listas será combinação dessas quatro jogadas. A filosofia do Kotlin é: se a lógica cabe em uma frase, jamais escreva cinco.

\booksection{Cartão de consulta rápida da sintaxe para quem parte do zero}

Não tente decorar toda a sintaxe num dia. Salve este cartão, consulte-o enquanto programa e, em duas semanas, a familiaridade vem naturalmente:

\begin{center}
\begin{tabularx}{\linewidth}{|X|X|}
\hline
\textbf{O que você quer fazer} & \textbf{Como se escreve} \tabularnewline
\hline
Definir variável somente leitura & \texttt{val x = 1} \tabularnewline
Definir variável mutável & \texttt{var x = 1} \tabularnewline
Definir função & fun nome(parâmetros): tipoDeRetorno \{ ... \} \tabularnewline
Decisão & if (condição) \{...\} else \{...\} \tabularnewline
Múltiplas escolhas & when \{ condição -> resultado \} \tabularnewline
Percorrer & for (x in lista) \{ ... \} \tabularnewline
Concatenação de strings & "Olá, \$name" \tabularnewline
Tratar nulos & \texttt{x?.\allowbreak length ?: 0} \tabularnewline
Definir data class & \texttt{data class X(val a: Int, val b: String)} \tabularnewline
\hline
\end{tabularx}
\end{center}

\begin{center}\small\color{gray}Cartão de consulta rápida da sintaxe Kotlin (edição iniciante)\end{center}

\begin{tcbpractice}{Mão na massa: jogue 10 minutos de Kotlin sem instalar nada}
Abra play.kotlinlang.org no navegador (o playground oficial online do Kotlin), cole e rode cada trecho de código deste capítulo, mude números e textos e observe como a saída muda. Esses 10 minutos de “brincar antes de estudar” vão deixar o início do Capítulo 51 duas vezes mais leve.

\end{tcbpractice}

\begin{tcbsummary}{Resumo do capítulo}
\begin{itemize}
  \item val primeiro, var com cautela; o tipo pode ser inferido automaticamente, ou declarado explicitamente depois dos dois-pontos
  \item Função = lógica empacotada; o template de string com \$ é a conveniência da casa
  \item if/when para decisões, for para laços; o when é o "coringa" do desenvolvimento Android
  \item A segurança contra nulos é o amuleto do Kotlin: tipos comuns recusam null; tipos anuláveis exigem tratamento com ? / ?:
  \item A data class é a caixa que guarda dados, um campo por linha, usada todos os dias em projetos reais
  \item 20\% da sintaxe fazem 80\% do trabalho: primeiro saiba usar, depois busque a maestria
\end{itemize}

\end{tcbsummary}
